\documentclass[11pt,openany]{book}

\usepackage[margin=1in]{geometry}
\usepackage{amsmath,amssymb,amsfonts,mathtools}
\usepackage{graphicx}
\usepackage{hyperref}
\usepackage{cleveref}
\usepackage{booktabs}
\usepackage{enumitem}
\usepackage{tikz}
\usetikzlibrary{quantikz}
\usepackage{braket}
\usepackage{microtype}
\usepackage{xcolor}

\hypersetup{
  colorlinks=true,
  linkcolor=blue!60!black,
  urlcolor=cyan!40!black,
  citecolor=green!40!black
}

\newcommand{\iqc}{Introduction to Quantum Computing}
\newcommand{\cutoff}{September 2026}

\title{\iqc\\[0.5em]\large From Classical Bits to Quantum Algorithms\\[1em]\normalsize 2026 Edition}
\author{An Independent Modern Textbook\\Research cutoff: \cutoff}
\date{}

\begin{document}
\frontmatter
\maketitle

\begin{abstract}
This textbook provides a mathematically careful yet accessible introduction to quantum computation and quantum information for readers who know basic algebra, trigonometry, and complex numbers but need not already know linear algebra, quantum mechanics, or quantum programming. It develops classical information and computation, reversible logic, single- and multi-qubit quantum theory, measurement, the Bloch sphere, linear-algebraic tools, entanglement, quantum circuits, error correction, protocols, and foundational quantum algorithms through Shor's factoring algorithm. A companion interactive website synchronizes notation, examples, and laboratory exercises with this PDF. Material labeled with the research cutoff reflects the state of the field as of \cutoff{} and should be re-checked as hardware and software evolve.
\end{abstract}

\tableofcontents

\chapter*{How to Use This Book}
\addcontentsline{toc}{chapter}{How to Use This Book}

Each chapter in the PDF corresponds to a section of the interactive web textbook at the project repository. When you see a \textbf{Interactive Lab} box, open the linked lab in your browser to manipulate states, run simulations, and verify calculations step by step.

\mainmatter

\chapter{Classical Information and Computation}

Physical systems can encode information by occupying distinguishable states. A classical bit is the simplest unit: it takes one of two values, $0$ or $1$. For $n$ bits there are $2^n$ distinct register states.

\section{States and Information}

If each bit is independent, an $n$-bit register has
\[
N = 2^n
\]
possible configurations. For $n=3$, the eight strings are $000,001,\ldots,111$. Bits may represent numbers, text, instructions, or any digital data once an encoding convention is fixed.

\begin{center}
\fbox{\parbox{0.9\textwidth}{\textbf{Interactive Lab 1:} Binary State Explorer --- select $n$ and inspect $2^n$ on the website.}}
\end{center}

\section{Binary Numbers}

Positional notation interprets a bit string $b_{n-1}\cdots b_0$ as
\[
\sum_{j=0}^{n-1} b_j 2^j.
\]
For example, $1011_2 = 1\cdot 8 + 0\cdot 4 + 1\cdot 2 + 1 = 11_{10}$.

\section{Logic Gates}

Boolean variables $A,B\in\{0,1\}$ combine through gates such as NOT ($\overline A$), AND ($A\land B$), OR ($A\lor B$), and XOR ($A\oplus B$). XOR satisfies $1\oplus 1 = 0$ and is its own inverse.

\section{Boolean Algebra}

Identities include $A\land 1 = A$, $A\lor 0 = A$, and $A\land\overline A = 0$. De Morgan's laws state
\[
\overline{A\land B} = \overline A \lor \overline B,\qquad
\overline{A\lor B} = \overline A \land \overline B.
\]

\section{Adders}

A half adder outputs sum $S = A\oplus B$ and carry $C = A\land B$. A full adder adds an incoming carry $C_{\mathrm{in}}$:
\[
S = A\oplus B\oplus C_{\mathrm{in}},\qquad
C_{\mathrm{out}} = (A\land B)\lor\bigl(C_{\mathrm{in}}\land(A\oplus B)\bigr).
\]
Ripple-carry addition propagates carries sequentially, introducing a depth cost that foreshadows circuit depth in quantum computing.

\section{Reversible Computation}

Reversible operations preserve enough information to recover inputs. NOT is reversible; AND is not. The Toffoli gate $(a,b,c)\mapsto(a,b,c\oplus ab)$ is reversible and can emulate AND when $c=0$. Quantum evolution before measurement must be reversible (unitary).

\section{Classical Error Correction}

Three-bit repetition encoding uses $0_L=000$ and $1_L=111$. A single bit flip is corrected by majority vote. Parity $p=a\oplus b$ summarizes two-bit information.

\section{Computational Complexity}

Growth classes include $O(1)$, $O(\log n)$, $O(n)$, $O(n\log n)$, $O(n^2)$, and $O(2^n)$. The class \textbf{P} contains problems solvable in polynomial time on a deterministic classical computer. \textbf{NP} is defined decision-theoretically; quantum computers do not automatically solve NP-complete problems in polynomial time.

\section{Turing Machines}

A Turing machine reads and writes symbols on a tape, changes internal state, and moves its head according to transition rules. The Church--Turing thesis asserts that realistic classical computation is captured by this model. Quantum computation asks whether every physically realizable process is efficiently simulable classically; this motivates the class \textbf{BQP}.

\section*{Exercises}
\begin{enumerate}[label=\arabic*.]
\item How many states does a 10-bit register have?
\item Simplify $\overline{A\land(B\lor C)}$ using De Morgan.
\item Add $1101_2 + 1011_2$ with ripple carry.
\end{enumerate}

\chapter{The Qubit}

A qubit generalizes a classical bit. Whereas a classical bit is either $0$ or $1$, a qubit may occupy a \emph{superposition}---a normalized linear combination of $|0\rangle$ and $|1\rangle$.

\section{What Is a Qubit?}

A pure one-qubit state is
\[
\ket{\psi} = \alpha\ket{0} + \beta\ket{1},\qquad \alpha,\beta\in\mathbb{C},\qquad |\alpha|^2+|\beta|^2=1.
\]
In the computational (Z) basis,
\[
\ket{0}=\begin{bmatrix}1\\0\end{bmatrix},\qquad
\ket{1}=\begin{bmatrix}0\\1\end{bmatrix}.
\]
A qubit is \emph{not} simply ``both 0 and 1 at once'' in a classical sense; amplitudes carry magnitude and phase, and interference governs observable statistics.

\section{Complex Amplitudes}

For $z=a+bi$, the conjugate is $z^*=a-bi$, magnitude $|z|=\sqrt{a^2+b^2}$, and polar form $z=re^{i\phi}$ with Euler's formula $e^{i\phi}=\cos\phi+i\sin\phi$.

\section{Measurement in the Z Basis}

Measuring $\ket{\psi}=\alpha\ket{0}+\beta\ket{1}$ in the Z basis yields outcome $0$ with probability $|\alpha|^2$ and outcome $1$ with probability $|\beta|^2$. After outcome $0$, the state collapses to $\ket{0}$.

\paragraph{Example.} For $\ket{\psi}=\frac{\sqrt3}{2}\ket{0}+\frac12\ket{1}$, we have $P(0)=3/4$ and $P(1)=1/4$.

\section{Other Measurement Bases}

The X basis uses $\ket{+}=\frac{\ket{0}+\ket{1}}{\sqrt2}$ and $\ket{-}=\frac{\ket{0}-\ket{1}}{\sqrt2}$. The Y basis uses $\ket{i}=\frac{\ket{0}+i\ket{1}}{\sqrt2}$ and $\ket{-i}=\frac{\ket{0}-i\ket{1}}{\sqrt2}$. Rewriting $\ket{\psi}$ in a chosen basis and squaring coefficient magnitudes gives measurement probabilities.

\section{Global and Relative Phase}

Multiplying $\ket{\psi}$ by a global phase $e^{i\gamma}$ does not change Z-basis probabilities. Relative phase between $\ket{0}$ and $\ket{1}$ \emph{does} matter: $\ket{+}$ and $\ket{-}$ have identical Z statistics but differ in X statistics.

\section{The Bloch Sphere}

Ignoring global phase, any pure qubit satisfies
\[
\ket{\psi}=\cos\frac{\theta}{2}\ket{0}+e^{i\phi}\sin\frac{\theta}{2}\ket{1},
\]
with Bloch coordinates $x=\sin\theta\cos\phi$, $y=\sin\theta\sin\phi$, $z=\cos\theta$. Important states include $\ket{0}\mapsto(0,0,1)$, $\ket{1}\mapsto(0,0,-1)$, $\ket{+}\mapsto(1,0,0)$, and $\ket{i}\mapsto(0,1,0)$.

\section{One-Qubit Gates}

Common gates include identity $I$, Pauli $X,Y,Z$, Hadamard $H$, phase $S$, and $T$. For example,
\[
X=\begin{bmatrix}0&1\\1&0\end{bmatrix},\quad
H=\frac{1}{\sqrt2}\begin{bmatrix}1&1\\1&-1\end{bmatrix},\quad
H\ket{0}=\ket{+}.
\]
Identities include $X^2=Y^2=Z^2=H^2=I$ and $S^2=Z$, $T^2=S$.

\begin{center}
\fbox{\parbox{0.9\textwidth}{\textbf{Interactive Labs 4--8:} Complex plane, measurement histograms, Bloch sphere, and gate explorer on the website.}}
\end{center}

\section*{Exercises}
\begin{enumerate}[label=\arabic*.]
\item Normalize $(2,1)^\top$ and compute $P(0)$.
\item Express $\ket{-}$ in the Z basis and compute $P(-)$ for $\ket{0}$.
\item Apply $H$ then $Z$ to $\ket{+}$ and give Bloch coordinates.
\end{enumerate}

\chapter{Linear Algebra for Quantum Computing}
\label{ch:linear-algebra}

Quantum states are vectors; quantum gates are matrices. This chapter develops the Dirac \emph{bra-ket} notation, inner products, projectors, outer products, completeness relations, and unitary matrices---the linear-algebraic language used in every subsequent chapter.

\section{Kets and Bras}
\label{sec:kets-bras}

\subsection{Column vectors and row covectors}

A ket $\ket{\psi}$ denotes a column vector in $\mathbb{C}^d$. For one qubit,
\begin{equation}
\ket{\psi}=\begin{bmatrix}\alpha\\\beta\end{bmatrix},\qquad
\bra{\psi}=\begin{bmatrix}\alpha^*&\beta^*\end{bmatrix}.
\label{eq:ket-bra}
\end{equation}
The dagger (conjugate transpose) maps kets to bras: $(\ket{\psi})^\dagger=\bra{\psi}$.

\subsection{Dual space}

Bras live in the dual space---the space of linear functionals on kets. The pairing $\bra{\phi}$ acting on $\ket{\psi}$ yields a complex number, defined next.

\section{Inner Products}
\label{sec:inner-products}

\subsection{Definition and properties}

For $\ket{\phi}=\begin{bmatrix}a\\b\end{bmatrix}$ and $\ket{\psi}=\begin{bmatrix}c\\d\end{bmatrix}$,
\begin{equation}
\braket{\phi|\psi}=a^*c+b^*d,\qquad
\|\ket{\psi}\|=\sqrt{\braket{\psi|\psi}}.
\label{eq:inner-product}
\end{equation}
Properties: conjugate symmetry $\braket{\phi|\psi}=\braket{\psi|\phi}^*$, linearity in the second slot, and positive definiteness $\braket{\psi|\psi}\ge 0$ with equality iff $\ket{\psi}=0$.

\subsection{Orthogonality and orthonormal bases}

Vectors are \emph{orthogonal} if $\braket{\phi|\psi}=0$. An \emph{orthonormal basis} satisfies $\braket{i|j}=\delta_{ij}$. The computational basis~\eqref{eq:computational-basis} is orthonormal.

\begin{example}
Compute $\braket{+|-}$ using~\eqref{eq:x-basis}:
\[
\braket{+|-}=\frac{1}{2}(\bra{0}+\bra{1})(\ket{0}-\ket{1})=\frac{1}{2}(1-1)=0.
\]
Thus $\ket{+}$ and $\ket{-}$ are orthogonal despite sharing the same $Z$-measurement statistics.
\end{example}

\section{Projection and Measurement}
\label{sec:projection}

\subsection{Born rule revisited}

Measuring in basis $\{\ket{m}\}$, the probability of outcome $m$ is
\begin{equation}
P(m)=|\braket{m|\psi}|^2.
\label{eq:born-rule}
\end{equation}
The amplitude $\braket{m|\psi}$ is sometimes called the \emph{overlap} of $\ket{\psi}$ with $\ket{m}$.

\subsection{Projectors}

The rank-one projector onto $\ket{m}$ is
\begin{equation}
P_m=\ket{m}\bra{m}.
\label{eq:projector}
\end{equation}
Then $P_m^2=P_m$, $P_m^\dagger=P_m$, and $P(m)=\bra{\psi}P_m\ket{\psi}$. For non-orthogonal measurements (POVMs), generalize to positive operators summing to identity---beyond our scope here.

\begin{example}
For $\ket{\psi}=\alpha\ket{0}+\beta\ket{1}$,
\[
P_0=\ket{0}\bra{0}=\begin{bmatrix}1&0\\0&0\end{bmatrix},\quad
P(0)=\bra{\psi}P_0\ket{\psi}=|\alpha|^2.
\]
\end{example}

\section{Outer Products and Completeness}
\label{sec:outer-completeness}

\subsection{Rank-one operators}

The outer product $\ket{\psi}\bra{\phi}$ is a $d\times d$ matrix. Example:
\begin{equation}
\ket{1}\bra{0}=\begin{bmatrix}0&0\\1&0\end{bmatrix}.
\label{eq:outer-product}
\end{equation}
Such operators appear in quantum channels and density matrices.

\subsection{Resolution of the identity}

Orthonormal basis $\{\ket{i}\}$ satisfies \emph{completeness}:
\begin{equation}
\sum_i \ket{i}\bra{i}=I.
\label{eq:completeness}
\end{equation}
In the computational basis, $\ket{0}\bra{0}+\ket{1}\bra{1}=I$. Any state expands as $\ket{\psi}=\sum_i \ket{i}\braket{i|\psi}$.

\begin{example}
Expand $\ket{+}$ in the $Z$ basis:
\[
\ket{+}=\ket{0}\braket{0|+}+\ket{1}\braket{1|+}=\frac{1}{\sqrt{2}}(\ket{0}+\ket{1}).
\]
\end{example}

\section{Unitary Matrices}
\label{sec:unitary}

\subsection{Definition}

A matrix $U$ is \emph{unitary} if
\begin{equation}
U^\dagger U = UU^\dagger = I,\qquad U^{-1}=U^\dagger.
\label{eq:unitary-def}
\end{equation}
Unitary matrices preserve inner products: $\braket{U\phi|U\psi}=\braket{\phi|\psi}$.

\subsection{Physical interpretation}

Quantum gates are unitary evolutions (before measurement). Normalization is preserved:
\begin{equation}
\braket{\psi'|\psi'}=\braket{\psi|U^\dagger U|\psi}=\braket{\psi|\psi}.
\label{eq:norm-preservation}
\end{equation}
Determinant has unit magnitude: $|\det U|=1$. Eigenvalues lie on the unit circle.

\begin{labbox}{9}{Matrix \& Unitary Checker}
Enter a $2\times 2$ matrix at \texttt{/playground/unitary-checker}. Verify $H^\dagger H=I$ and confirm that a random complex matrix fails unitarity unless specially constructed.
\end{labbox}

\begin{example}
Show $H$ is unitary by computing $H^\dagger H$:
\[
H^\dagger H=\frac{1}{2}\begin{bmatrix}1&1\\1&-1\end{bmatrix}\begin{bmatrix}1&1\\1&-1\end{bmatrix}
=\frac{1}{2}\begin{bmatrix}2&0\\0&2\end{bmatrix}=I.
\]
\end{example}

\subsection{Special unitary group}

Matrices with $\det U=1$ form $\mathrm{SU}(2)$, the double cover of rotations on the Bloch sphere. Pauli matrices $X,Y,Z$ generate $\mathrm{SU}(2)$ up to global phase.

\section{Linear Operators on Composite Systems (Preview)}
\label{sec:tensor-preview}

When subsystems combine, operators act as tensor products: $A\otimes B$ applies $A$ to the first subsystem and $B$ to the second. Chapter~\ref{ch:multi} develops this systematically. For now, note that unitarity extends: $(A\otimes B)^\dagger=A^\dagger\otimes B^\dagger$.

\section*{Exercises}
\begin{enumerate}[label=\arabic*.]
\item Compute $\braket{0|1}$, $\braket{+|\psi}$ for $\ket{\psi}=\ket{0}$, and $\braket{i|+i}$ where $\ket{i}=\ket{+i}$ in the $Y$ basis.
\item Verify~\eqref{eq:completeness} explicitly for the $X$ basis states $\ket{\pm}$.
\item Write $\ket{1}\bra{0}$ and $\ket{+}\bra{-}$ as $2\times 2$ matrices.
\item Show that if $U$ is unitary, so is $U^\dagger$.
\item Prove $\|U\ket{\psi}\|=\|\ket{\psi}\|$ using~\eqref{eq:unitary-def}.
\item For projector $P_+=\ket{+}\bra{+}$, compute $\mathrm{Tr}(P_+)$ and $P_+^2$.
\item Express $X$ as a linear combination of outer products in the $Z$ basis.
\item Why must quantum gate matrices be unitary rather than merely invertible?
\end{enumerate}

\chapter{Composite Quantum Systems}

Multiple qubits are described by tensor products. An $n$-qubit computational basis has $2^n$ states.

\section{Tensor Products}

Two-qubit product states satisfy $\ket{\psi}\otimes\ket{\phi}$; for example $\ket{0}\otimes\ket{1}=\ket{01}$. A general two-qubit state is
\[
\ket{\psi}=c_{00}\ket{00}+c_{01}\ket{01}+c_{10}\ket{10}+c_{11}\ket{11},
\]
with $\sum|c_{ij}|^2=1$.

\section{Entanglement}

A state is \emph{separable} if it equals a product $\ket{\alpha}\otimes\ket{\beta}$; otherwise it is \emph{entangled}. Bell states include
\[
\ket{\Phi^+}=\frac{\ket{00}+\ket{11}}{\sqrt2}.
\]
Measuring one qubit of $\ket{\Phi^+}$ in the Z basis yields correlated outcomes; this does not enable faster-than-light signaling.

\section{Multi-Qubit Gates}

The CNOT gate acts as $\ket{a,b}\mapsto\ket{a,a\oplus b}$. Creating $\ket{\Phi^+}$ from $\ket{00}$ uses $H$ on the first qubit followed by CNOT. SWAP exchanges qubits; Toffoli generalizes classical reversible AND.

\section{No-Cloning Theorem}

No unitary process can clone an unknown arbitrary qubit: cloning would require $\braket{\psi|\phi}=\braket{\psi|\phi}^2$ for all pairs, which fails unless states are orthogonal.

\section{Universal Gate Sets}

The set $\{H,T,\mathrm{CNOT}\}$ is universal: any $n$-qubit unitary can be approximated to arbitrary precision using finitely many gates from a universal set (Solovay--Kitaev theorem, stated without proof).

\begin{center}
\fbox{\parbox{0.9\textwidth}{\textbf{Interactive Labs 10--12:} Tensor products, Bell pairs, entanglement measurement.}}
\end{center}

\section*{Exercises}
\begin{enumerate}[label=\arabic*.]
\item Write $\ket{+}\otimes\ket{0}$ in the four-bit computational basis.
\item Show $\ket{\Phi^+}$ is not a product state.
\item Apply CNOT to $\frac{\ket{00}+\ket{10}}{\sqrt2}$.
\end{enumerate}

\chapter{Quantum Error Correction}

Real devices decohere. Error correction protects \emph{logical} qubits using redundancy in \emph{physical} qubits.

\section{Noise and Decoherence}

Errors include bit flips ($X$), phase flips ($Z$), depolarizing noise, relaxation ($T_1$), and dephasing ($T_2$). As of the \cutoff{} research cutoff, large-scale fault-tolerant quantum computation remains an active engineering goal; recent experiments demonstrate improving logical error rates with surface-code scaling \cite{acharya2024}, while commercial roadmaps should be distinguished from established theorems.

\section{Three-Qubit Bit-Flip Code}

Logical encoding: $\ket{0_L}=\ket{000}$, $\ket{1_L}=\ket{111}$. An arbitrary logical qubit $\alpha\ket{0_L}+\beta\ket{1_L}$ can suffer a single $X$ error on one wire; parity measurements reveal the syndrome without collapsing $\alpha,\beta$.

\section{Phase-Flip Code}

Phase errors flip $\ket{+}\leftrightarrow\ket{-}$. Conjugation by $H$ exchanges $X$ and $Z$, reducing phase-flip correction to bit-flip correction in the rotated basis.

\section{Shor Nine-Qubit Code}

The Shor code combines bit- and phase-flip protection using nine physical qubits per logical qubit. Modern research emphasizes surface codes and other topological schemes with higher thresholds and more hardware-friendly layouts.

\section{Error Mitigation vs.\ Correction}

In the NISQ era \cite{preskill2018}, error \emph{mitigation} techniques (zero-noise extrapolation, probabilistic error cancellation) reduce bias without full logical encoding. Mitigation is not a substitute for fault tolerance at scale.

\begin{center}
\fbox{\parbox{0.9\textwidth}{\textbf{Interactive Lab 14:} Inject bit-flip errors and correct with syndromes.}}
\end{center}

\section*{Exercises}
\begin{enumerate}[label=\arabic*.]
\item What syndrome indicates an $X$ error on qubit 2 in the three-bit code?
\item Why cannot you clone $\ket{\psi}$ to protect it?
\item Distinguish physical and logical qubits in one paragraph.
\end{enumerate}

\chapter{Quantum Circuits and Programming}
\label{ch:circuits}

Quantum algorithms are expressed as \emph{circuits}: wires represent qubits, time flows left to right, boxes are gates, control dots mark controlled operations, and measurement symbols denote readout. This chapter formalizes the circuit model, walks through Bell-state synthesis, surveys gate sets and OpenQASM, and discusses depth and resource costs---connecting mathematics to the software tools used on real hardware.

\section{The Circuit Model}
\label{sec:circuit-model}

\subsection{Unitary evolution}

An $n$-qubit circuit transforms $\ket{\psi_{\mathrm{in}}}$ to
\begin{equation}
\ket{\psi_{\mathrm{out}}}=U_k\cdots U_2 U_1\ket{\psi_{\mathrm{in}}},
\label{eq:circuit-evolution}
\end{equation}
where each $U_i$ acts unitarily on some subset of qubits (often tensored with identity on others). Measurement at the end samples classical outcomes from Born probabilities. Circuits are \emph{acyclic}: no feedback unless classical control is explicitly modeled.

\subsection{Diagrammatic conventions}

Wires index qubits from top to bottom; time flows left to right. A solid control dot on one wire with $\oplus$ on another denotes CNOT. Single-qubit boxes label gates ($H$, $T$, etc.). Double lines occasionally denote classical bits conditioned on measurement outcomes.

\subsection{Depth, width, and gate count}

\begin{itemize}[noitemsep]
\item \textbf{Width:} number of qubits $n$.
\item \textbf{Gate count:} total gates applied.
\item \textbf{Depth:} number of sequential \emph{layers} of gates that can run in parallel on disjoint qubits.
\end{itemize}
Depth determines latency on hardware with sequential gate calibration; gate count relates to total error budget. On superconducting devices, two-qubit gates along the same coupling edge often cannot execute simultaneously.

\begin{example}
Two CNOTs sharing a control qubit cannot run in parallel if they use the same control line in conflicting ways; two single-qubit gates on different qubits \emph{can} occupy one depth layer. A circuit with $n$ sequential CNOTs on a chain has depth $\Omega(n)$ even if gate count is also $n$.
\end{example}

\subsection{Measurement and classical control}

Mid-circuit measurement feeds classical registers that condition later gates (\emph{dynamic circuits}). This extends~\eqref{eq:circuit-evolution} to conditional operations $U_k(c_1,\ldots,c_m)$ where $c_j$ are prior measurement outcomes---essential for teleportation and error correction.

\section{Building Bell States}
\label{sec:bell-circuit}

\subsection{Standard construction}

The canonical Bell circuit prepares $\ket{\Phi^+}$ from $\ket{00}$:
\begin{equation}
\ket{00}\xrightarrow{H_0}\xrightarrow{\mathrm{CNOT}_{0,1}}\ket{\Phi^+}
=\frac{\ket{00}+\ket{11}}{\sqrt{2}}.
\label{eq:bell-circuit}
\end{equation}
Here $H_0$ denotes Hadamard on qubit~0; $\mathrm{CNOT}_{0,1}$ uses qubit~0 as control and qubit~1 as target.

\subsection{Step-by-step state update}

\begin{align*}
\ket{00} &\xrightarrow{H\otimes I} \frac{1}{\sqrt{2}}(\ket{00}+\ket{10}), \\
&\xrightarrow{\mathrm{CNOT}} \frac{1}{\sqrt{2}}(\ket{00}+\ket{11})=\ket{\Phi^+}.
\end{align*}
Other Bell states arise from additional single-qubit Pauli rotations before or after this core.

\begin{labbox}{CB}{Circuit Builder}
Open \texttt{/playground/circuit-builder}. Drag $H$ and CNOT onto two qubits, simulate step-by-step, and read the statevector. Confirm amplitudes match~\eqref{eq:bell-states}.
\end{labbox}

\section{Gate Sets and Native Hardware}
\label{sec:gate-sets}

\subsection{Universal versus native gates}

Theory uses $\{H,T,\mathrm{CNOT}\}$; hardware exposes \emph{native} gates---e.g.\ $R_x,R_y,R_z$ and CNOT on fixed connectivity graphs. \emph{Compilation} transpiles abstract circuits into native pulses with SWAP routing for limited connectivity.

\subsection{Clifford+T and T-count}

Fault-tolerant architectures often separate fast \emph{Clifford} gates ($H,S,\mathrm{CNOT}$) from costly \emph{T} gates, which require magic-state distillation. $T$-count dominates resource estimates for Shor's algorithm.

\section{OpenQASM}
\label{sec:openqasm}

\subsection{Language overview}

Open Quantum Assembly Language (OpenQASM) describes circuits textually. OpenQASM~3.0 adds classical control flow, gate modifiers, and richer typing. A minimal Bell preparation:

\begin{verbatim}
OPENQASM 3.0;
include "stdgates.inc";
qubit[2] q;
h q[0];
cx q[0], q[1];
\end{verbatim}

\subsection{Registers and timing}

OpenQASM declares \texttt{qubit[n]} quantum and \texttt{bit[m]} classical registers. Statements execute in order unless grouped into \texttt{box} blocks for timing semantics. Gate arguments specify qubit indices; creg stores measurement outcomes for conditional gates such as \texttt{if (c == 1) x q[0];}.

\subsection{Reading and writing QASM}

Export from visual builders or SDKs (Qiskit \texttt{qasm3.dumps}, Cirq \texttt{to\_qasm}). Always map each statement back to unitaries on the Hilbert space; \texttt{cx} is CNOT, \texttt{h} is Hadamard, \texttt{rz(theta)} is $R_z(\theta)$.

\begin{example}
The OpenQASM snippet above corresponds to $U=\mathrm{CNOT}_{0,1}\,(H\otimes I)$ applied left-to-right as $U\ket{00}=\ket{\Phi^+}$. Reversing gate order in the file changes the unitary product accordingly.
\end{example}

\begin{example}
The unitary for $H\otimes I$ on two qubits is
\[
H\otimes I=\frac{1}{\sqrt{2}}\begin{bmatrix}
1&1&0&0\\1&-1&0&0\\0&0&1&1\\0&0&1&-1
\end{bmatrix}.
\]
Apply it to $\ket{00}$ before CNOT to derive~\eqref{eq:bell-circuit}.
\end{example}

\section{Depth and Resource Analysis}
\label{sec:depth}

\subsection{Ripple adders revisited}

An $n$-bit ripple-carry adder built from Toffoli and CNOT has depth $O(n)$ in the worst case---analogous classical depth from Chapter~\ref{ch:classical}. Quantum carry-lookahead and Fourier adders reduce depth at the cost of ancillas.

\subsection{Simulation cost}

Exact state-vector simulation of $n$ qubits requires $O(2^n)$ complex amplitudes. Clifford circuits (generated by $\{H,S,\mathrm{CNOT}\}$ without $T$) simulate in polynomial time (Gottesman--Knill), explaining why random-circuit supremacy demos use non-Clifford gates \cite{arute2019}.

\section{Modern Software Ecosystem (2026)}
\label{sec:software}

Production quantum software typically combines:
\begin{itemize}[noitemsep]
\item \textbf{Circuit construction:} OpenQASM~3, Qiskit, Cirq, PennyLane, and vendor SDKs.
\item \textbf{Simulation:} local state-vector simulators (exponential in $n$) and stabilizer simulators for Clifford circuits.
\item \textbf{Execution:} cloud access to superconducting, trapped-ion, neutral-atom, and photonic backends with calibration data and error models.
\end{itemize}

Map API calls (\texttt{h(0)}, \texttt{cx(0,1)}) to unitaries and tensor structure before trusting black-box outputs.

\section*{Exercises}
\begin{enumerate}[label=\arabic*.]
\item Write the $4\times 4$ matrix for $H\otimes I$ and verify $ (H\otimes I)\ket{00}$.
\item Count depth (qualitatively) of an $n$-bit ripple-carry adder from Toffoli gates.
\item Export a Bell circuit to OpenQASM: which gates appear?
\item Why does CNOT not equal $\mathrm{CNOT}_{1,0}$? Give the matrix of swapped control/target.
\item Estimate memory for simulating $n=30$ qubits in double precision (bytes $\approx 16\cdot 2^n$).
\item Which gates in $\{H,S,\mathrm{CNOT}\}$ generate the Clifford group?
\item Transpile $T$ into $R_z(\pi/4)$ up to global phase.
\item Explain why measurement is not modeled as a unitary gate in~\eqref{eq:circuit-evolution}.
\end{enumerate}

\chapter{Entanglement and Quantum Protocols}
\label{ch:protocols}

Entanglement enables communication and cryptography protocols impossible with classical bits alone, while respecting the \emph{no-signaling} principle: local measurements cannot transmit controllable faster-than-light messages. This chapter covers EPR paradox and hidden variables, the CHSH inequality, superdense coding, quantum teleportation, and BB84 quantum key distribution.

\section{EPR and Hidden Variables}
\label{sec:epr}

\subsection{The EPR argument}

Einstein, Podolsky, and Rosen (1935) highlighted that entangled pairs exhibit perfect correlations when both parties measure in the same basis, yet each qubit alone appears maximally random. They questioned whether quantum mechanics is \emph{complete} or whether \emph{local hidden variables} (LHV) predetermine outcomes.

\subsection{Local realism versus quantum predictions}

An LHV model assigns definite values to observables before measurement, with outcomes determined locally without faster-than-light influence. Bell and CHSH showed certain correlations violate inequalities any LHV model must satisfy.

\begin{example}
Shared $\ket{\Phi^+}$: Alice and Bob both measure $Z$. Outcomes are perfectly correlated ($00$ or $11$) yet individually random with $P(0)=P(1)=1/2$. Classically, one might imagine hidden bits copied at source; Bell tests rule out \emph{local} hidden variables for maximal entanglement.
\end{example}

\section{CHSH Inequality}
\label{sec:chsh}

\subsection{Setup}

Alice chooses measurement settings $a,a'\in\{0,1\}$; Bob chooses $b,b'$. Define correlation
\begin{equation}
E(a,b)=P(00)+P(11)-P(01)-P(10)
\label{eq:chsh-correlation}
\end{equation}
for $\pm1$ outcomes mapped to $\{+1,-1\}$.

\subsection{Classical and quantum bounds}

The CHSH combination
\begin{equation}
S=E(a,b)+E(a,b')+E(a',b)-E(a',b')
\label{eq:chsh-s}
\end{equation}
satisfies $|S|\le 2$ for any local hidden-variable theory. Quantum mechanics permits $|S|\le 2\sqrt{2}$ (Tsirelson bound), achieved by suitable measurements on $\ket{\Phi^-}$ or singlet states.

\begin{labbox}{14}{Bell/CHSH Experiment}
Simulate \texttt{/playground/chsh}: choose measurement angles, estimate $E(a,b)$ from shot statistics, and compute $S$. Push toward the quantum limit $2\sqrt{2}$.
\end{labbox}

\begin{example}
For singlet $\ket{\Psi^-}=\frac{1}{\sqrt{2}}(\ket{01}-\ket{10})$, opposite measurements in the Bloch plane can yield $E=-\cos(\theta_A-\theta_B)$. Optimizing settings gives $|S|=2\sqrt{2}$.
\end{example}

\section{Superdense Coding}
\label{sec:superdense}

\subsection{Protocol}

Alice and Bob share $\ket{\Phi^+}_{AB}$. Alice encodes two classical bits into one qubit by applying one of $\{I,X,Z,XZ\}$ to her half, then sends her qubit to Bob. Bob performs a Bell-basis measurement decoding $(00),(01),(10),(11)$.

\subsection{Resource trade-off}

Two classical bits travel on one qubit \emph{plus} one shared ebit (entang bit). Entanglement is consumable but not cloneable; pre-shared entanglement is a resource.

\begin{labbox}{15}{Superdense Coding}
Walk through \texttt{/playground/superdense}: pick a two-bit message, apply Alice's encoding, and verify Bob's Bell measurement recovers the message with certainty.
\end{labbox}

\begin{example}
Alice sends bit pair $(1,0)$ using $Z$. On $\ket{\Phi^+}$, $Z_A\ket{\Phi^+}=\ket{\Phi^-}$. Bob's Bell measurement distinguishes $\ket{\Phi^+}$ from $\ket{\Phi^-}$, decoding $(1,0)$.
\end{example}

\section{Quantum Teleportation}
\label{sec:teleportation}

\subsection{Protocol steps}

Unknown state $\ket{\psi}=\alpha\ket{0}+\beta\ket{1}$ on qubit~0 teleports to Bob using shared $\ket{\Phi^+}_{12}$ on qubits~1 (Alice) and~2 (Bob):
\begin{enumerate}[noitemsep]
\item Alice applies CNOT$_{0\to 1}$ then $H_0$.
\item Alice measures qubits~0 and~1 in $Z$, obtaining two classical bits.
\item Bob applies corrections $X^{m_1}Z^{m_0}$ to qubit~2 based on outcomes $(m_0,m_1)$.
\end{enumerate}
Bob's qubit ends in $\ket{\psi}$; Alice's original is destroyed.

\subsection{No cloning and no signaling}

Teleportation transfers quantum state using one ebit and two classical bits. It does not copy $\ket{\psi}$ (no-cloning) and cannot signal faster than light because Bob's state is mixed until classical bits arrive.

\begin{labbox}{16}{Teleportation Simulator}
Run \texttt{/playground/teleportation} with random $\ket{\psi}$. Track Alice's measurements and Bob's conditional corrections.
\end{labbox}

\begin{example}
After Alice measures $\ket{\Phi^+}$ shared pair and gets $(m_0,m_1)=(0,0)$, Bob's qubit is $X^0Z^0\ket{\psi}=\ket{\psi}$ without correction if the circuit normalization matches standard conventions; other outcomes require $X$ and/or $Z$.
\end{example}

\section{BB84 Quantum Key Distribution}
\label{sec:bb84}

\subsection{Protocol outline}

\begin{enumerate}[noitemsep]
\item Alice chooses random bits and random bases ($Z$ or $X$), encodes, and sends qubits.
\item Bob measures in random bases.
\item They publicly compare bases (not outcomes); matching bases yield sifted key bits.
\item They sample a subset to estimate error rate; elevated errors suggest eavesdropping.
\end{enumerate}

\subsection{Security intuition}

An eavesdropper Eve cannot measure unknown states without disturbance (no-cloning, measurement back-action). Intercept-resend attacks increase bit-error rates detectable in the sifted key.

\begin{labbox}{17}{BB84 Simulator}
Use \texttt{/playground/bb84} with and without an eavesdropper. Compare error rates on the sifted key.
\end{labbox}

\begin{example}
Alice encodes bit~1 in $X$ basis as $\ket{-}$. Bob measures in $Z$, obtaining random outcomes---sifted key uses only rounds where bases matched. If Eve intercepts with random bases, $\approx 25\%$ of sifted bits err for unbiased encoding.
\end{example}

\section*{Exercises}
\begin{enumerate}[label=\arabic*.]
\item After Alice measures her qubit of $\ket{\Phi^+}$ in $Z$ and gets $1$, what is Bob's qubit (before any correction)?
\item Why does Eve's intercept-resend attack increase the error rate in BB84?
\item Sketch how $S$ in~\eqref{eq:chsh-s} is computed from four correlation settings.
\item How many classical bits and ebits does teleportation consume per unknown qubit?
\item Can superdense coding work without pre-shared entanglement? Justify briefly.
\item Write the four Bell measurement projectors as $4\times 4$ matrices.
\item Explain no-signaling: why can't Alice's local choice of $Z$ versus $X$ measurement send a message to Bob instantly?
\item Compare BB84 security assumptions with RSA security assumptions in one paragraph.
\end{enumerate}

\chapter{Quantum Algorithms}
\label{ch:algorithms}

Quantum algorithms exploit interference and phase to outperform classical strategies for \emph{specific} problems. This chapter develops the oracle model, Deutsch and Deutsch--Jozsa algorithms, Bernstein--Vazirani, Simon's algorithm, Grover search, the quantum Fourier transform, phase estimation, and Shor's factoring algorithm---with complexity comparisons throughout.

\section{Oracles and Phase Kickback}
\label{sec:oracle}

\subsection{Standard oracle}

A Boolean function $f:\{0,1\}^n\to\{0,1\}$ is accessed via a unitary
\begin{equation}
U_f\ket{x}\ket{y}=\ket{x}\ket{y\oplus f(x)}.
\label{eq:standard-oracle}
\end{equation}
The second register is an \emph{ancilla}; XOR implements reversible function evaluation.

\subsection{Phase oracle via kickback}

Initialize the ancilla in $\ket{-}=\frac{1}{\sqrt{2}}(\ket{0}-\ket{1})$. Then
\begin{equation}
U_f\ket{x}\ket{-}=\ket{x}\frac{(-1)^{f(x)}\ket{0}-\ket{1}}{\sqrt{2}}=(-1)^{f(x)}\ket{x}\ket{-}.
\label{eq:phase-kickback}
\end{equation}
Ignoring the unchanged ancilla defines the \emph{phase oracle}
\begin{equation}
O_f\ket{x}=(-1)^{f(x)}\ket{x}.
\label{eq:phase-oracle}
\end{equation}
Phase kickback is the workhorse of many algorithms.

\begin{labbox}{18}{Deutsch Algorithm}
Implement \texttt{/playground/deutsch}: test constant versus balanced $f:\{0,1\}\to\{0,1\}$ in one query using phase kickback.
\end{labbox}

\begin{example}
For $f(x)=x$ (balanced), $O_f\ket{0}=+\ket{0}$ and $O_f\ket{1}=-\ket{1}$: the oracle acts as $Z$ on the query qubit---detectable after $H$ interference.
\end{example}

\section{Deutsch and Deutsch--Jozsa}
\label{sec:deutsch-jozsa}

\subsection{Deutsch's problem}

For $n=1$, distinguish \emph{constant} ($f(0)=f(1)$) from \emph{balanced} ($f(0)\neq f(1)$) functions. Classically, two queries may be needed in the worst case; Deutsch's algorithm uses one.

\subsection{Circuit outline}

Prepare $\ket{+}\ket{-}$, apply $O_f$, apply $H$ on the query qubit, measure. Constant $f$ yields $\ket{0}$ with certainty; balanced $f$ yields $\ket{1}$.

\subsection{Deutsch--Jozsa generalization}

For $f:\{0,1\}^n\to\{0,1\}$ promised constant or balanced, classical worst-case queries are $2^{n-1}+1$; Deutsch--Jozsa uses one query with $n$ qubits plus ancilla.

\begin{labbox}{19}{Deutsch--Jozsa}
Explore \texttt{/playground/deutsch-jozsa} for $n=2,3$. Verify constant functions never produce all-ones parity pattern on the query register.
\end{labbox}

\begin{example}
The four $n=1$ functions: $f_0(x)=0$, $f_1(x)=1$ (constant); $f_2(x)=x$, $f_3(x)=1-x$ (balanced). Deutsch distinguishes the two classes in one shot.
\end{example}

\section{Bernstein--Vazirani}
\label{sec:bv}

\subsection{Hidden string problem}

Given $f(x)=s\cdot x \bmod 2$ (inner product mod~2) for hidden $s\in\{0,1\}^n$, find $s$. Classically, $n$ queries are required (one per bit); Bernstein--Vazirani uses one query after preparing $\ket{+}^{\otimes n}$ and applying $O_f$.

\subsection{Query complexity}

The algorithm demonstrates an exponential separation in query complexity for this structured promise problem---a template for Simon and Shor.

\begin{labbox}{20}{Bernstein--Vazirani}
At \texttt{/playground/bernstein-vazirani}, choose random $s$ and recover it from measurement outcomes in the $Z$ basis after $H^{\otimes n}$.
\end{labbox}

\section{Simon's Algorithm}
\label{sec:simon}

\subsection{Periodicity promise}

For $f:\{0,1\}^n\to\{0,1\}^n$ with hidden period $s\neq 0$ such that $f(x)=f(x\oplus s)$, find $s$. Classical algorithms need exponential queries; Simon needs $O(n)$ quantum queries plus classical post-processing.

\subsection{Linear algebra over $\mathbb{F}_2$}

Measurement outcomes $y$ satisfy $y\cdot s = 0 \pmod 2$. Collecting $\approx n$ independent equations yields $s$ by Gaussian elimination over $\mathbb{F}_2$.

\begin{labbox}{21}{Simon's Algorithm}
Simulate \texttt{/playground/simon}: observe random $y$ strings orthogonal to $s$ and solve for $s$ classically.
\end{labbox}

\begin{example}
If $s=110$ and measured $y_1=101$, $y_2=011$, then $y_1\cdot s=1+0+0=1\neq 0$---discard dependent samples until independent constraints determine $s$.
\end{example}

\section{Grover Search}
\label{sec:grover}

\subsection{Unstructured search}

Search space size $N=2^n$; exactly one marked state $\ket{w}$. Classically, $\Theta(N)$ queries are needed in the worst case; Grover achieves $O(\sqrt{N})$ oracle queries \cite{grover1996}.

\subsection{Grover iterate}

Define oracle $O_w=I-2\ket{w}\bra{w}$ and diffusion operator $D=2\ket{s}\bra{s}-I$ where $\ket{s}=\frac{1}{\sqrt{N}}\sum_x\ket{x}$. The iterate $G=DO_w$ rotates the amplitude vector toward $\ket{w}$ by angle $\approx 2\arcsin(1/\sqrt{N})$ per application.

\subsection{Optimal iteration count}

Approximately
\begin{equation}
k\approx \left\lfloor\frac{\pi}{4}\sqrt{N}\right\rfloor
\label{eq:grover-iterations}
\end{equation}
iterations maximize success probability.

\begin{labbox}{22}{Grover Search}
Run \texttt{/playground/grover} with $N=8$, marked $\ket{011}$. Compare $k=2$ versus $k=5$ success rates with~\eqref{eq:grover-iterations}.
\end{labbox}

\begin{example}
For $N=8$, $\lfloor\frac{\pi}{4}\sqrt{8}\rfloor=\lfloor 2.22\rfloor=2$. After two Grover iterations starting from uniform superposition, $P(\ket{w})\approx 1$ for the unique marked state.
\end{example}

\section{Quantum Fourier Transform}
\label{sec:qft}

\subsection{Definition}

On $N=2^n$ basis states,
\begin{equation}
\mathrm{QFT}_N\ket{x}=\frac{1}{\sqrt{N}}\sum_{k=0}^{N-1} e^{2\pi i xk/N}\ket{k}.
\label{eq:qft}
\end{equation}
The QFT is unitary and efficiently implementable with $O(n^2)$ gates, but amplitudes are \emph{not} all readable at once---measurement samples one outcome.

\begin{labbox}{23}{QFT Visualizer}
Inspect \texttt{/playground/qft} for $n=3$: apply QFT to $\ket{001}$ and compare magnitude patterns before measurement.
\end{labbox}

\section{Quantum Phase Estimation}
\label{sec:qpe}

\subsection{Problem}

Given unitary $U$ and eigenstate $\ket{u}$ with $U\ket{u}=e^{2\pi i\phi}\ket{u}$, estimate $\phi$ to $m$ bits using $m$ ancilla qubits and controlled powers $U^{2^j}$.

\subsection{Binary fraction expansion}

Phase estimation implements the QFT on controlled-$U$ outcomes, yielding an $m$-bit approximation to $\phi$. Shor's algorithm uses QPE on modular exponentiation unitaries.

\begin{labbox}{24}{Phase Estimation}
Use \texttt{/playground/phase-estimation} with $U=R_z$ on an eigenstate and $m=3$. Compare readout to $\phi=1/3\approx 0.333\ldots$ in binary.
\end{labbox}

\begin{example}
With $m=3$, binary digits approximate $\phi=1/3=0.\overline{010}_2$. Rounding errors arise from finite $m$; repeat with larger $m$ for precision.
\end{example}

\section{Shor's Factoring Algorithm}
\label{sec:shor}

\subsection{Reduction to period finding}

Factoring $N$ reduces to finding period $r$ of $f(x)=a^x\bmod N$ for random $a$ coprime to $N$. If $r$ is even and $a^{r/2}\not\equiv -1\pmod N$, then $\gcd(a^{r/2}\pm 1,N)$ yields nontrivial factors.

\subsection{Example $N=15$, $a=2$}

The period is $r=4$ since $2^4=16\equiv 1\pmod{15}$. Then
\[
\gcd(2^2+1,15)=\gcd(5,15)=5,\qquad
\gcd(2^2-1,15)=\gcd(3,15)=3.
\]
Quantum speedup lies in finding $r$ via QPE; remaining steps are classical \cite{shor1994}.

\begin{labbox}{25}{Period Explorer}
Explore \texttt{/playground/period-finding} for modular exponentiation and period $r$ of $f(x)=a^x\bmod N$.
\end{labbox}

\begin{labbox}{26}{Shor Demo}
Complete the factoring pipeline at \texttt{/playground/shor} using QPE-derived periods and classical $\gcd$ steps.
\end{labbox}

\subsection{Complexity summary}

Table~\ref{tab:algo-complexity} collects query/time scalings covered in this chapter; see Appendix~C for a fuller reference.

\begin{table}[H]
\centering
\caption{Representative query complexity advantages (promise problems).}
\label{tab:algo-complexity}
\begin{tabular}{@{}lll@{}}
\toprule
Problem & Classical (worst case) & Quantum \\
\midrule
Deutsch ($n=1$) & 2 queries & 1 \\
Deutsch--Jozsa & $2^{n-1}+1$ & 1 \\
Bernstein--Vazirani & $n$ & 1 \\
Simon & $O(2^{n/2})$ & $O(n)$ \\
Grover & $O(N)$ & $O(\sqrt{N})$ \\
Factoring (Shor) & no known poly-time & poly-time (ideal) \\
\bottomrule
\end{tabular}
\end{table}

\section*{Exercises}
\begin{enumerate}[label=\arabic*.]
\item Which of the four $n=1$ Boolean functions are constant versus balanced?
\item Run Grover on $N=8$ with marked $\ket{011}$; how many iterations does~\eqref{eq:grover-iterations} predict?
\item For QPE with $m=3$, what binary fraction approximates $\phi=1/3$?
\item Construct a phase oracle for $f(x)=1$ (constant) and show Deutsch outputs $\ket{0}$.
\item Why does Grover not achieve exponential speedup for NP-complete search?
\item Write the QFT matrix for $N=4$ and verify unitarity.
\item For Simon, why are measured $y$ vectors uniform over $\{0,1\}^n$ before post-processing?
\item Outline Shor's classical post-processing after obtaining period $r$.
\end{enumerate}

\chapter{Modern Quantum Computing in 2026}
\label{ch:modern}

This chapter surveys the landscape a serious introductory student should know beyond foundational algorithms: hardware platforms, surface codes and logical qubits, the NISQ era, software ecosystems, fault-tolerance roadmaps, and post-quantum cryptography. Statements are tagged by evidence type; company roadmaps are not scientific facts and must be verified against peer-reviewed literature.

\section{Physical Qubit Platforms}
\label{sec:hardware}

\subsection{Leading modalities}

As of the \cutoff{} cutoff, several qubit modalities compete:
\begin{itemize}[noitemsep]
\item \textbf{Superconducting transmons:} microwave control, fast gates, planar fabrication; short $T_1,T_2$ relative to ions.
\item \textbf{Trapped ions:} long coherence, high fidelity, all-to-all connectivity in chains; slower gates.
\item \textbf{Neutral atoms:} reconfigurable arrays via optical tweezers; promising for error-correction layouts.
\item \textbf{Photonic:} room-temperature operation; challenges in deterministic two-qubit gates.
\item \textbf{Spin qubits:} semiconductor integration potential; scaling research active.
\end{itemize}

\subsection{Metrics that matter}

Compare platforms using gate fidelity ($>99.9\%$ two-qubit targets), coherence times ($T_1,T_2$), connectivity graph, and reproducible cross-calibrated benchmarks---not qubit count alone.

\begin{example}
A device advertising ``1000+ qubits'' with $10^{-2}$ two-qubit error rates cannot yet run fault-tolerant Shor at cryptographically relevant scales without extensive QEC overhead \cite{preskill2018}.
\end{example}

\section{Logical Qubits and the Surface Code}
\label{sec:surface-code}

\subsection{From physical to logical}

Fault-tolerant computation stores information in \emph{logical} qubits encoded across many \emph{physical} qubits. Syndrome extraction detects errors faster than they accumulate. The \emph{threshold theorem} guarantees arbitrary low logical error rates if physical errors per gate stay below a code-specific threshold.

\subsection{Surface code highlights}

The surface code \cite{acharya2024} arranges stabilizer checks on a 2D lattice with local interactions suitable for superconducting and neutral-atom layouts. Recent experiments demonstrate \emph{logical} error rates below \emph{physical} rates at increasing code distance---a milestone toward scalable QEC, though resource estimates for chemistry or RSA-scale factoring remain substantial.

\begin{example}
Distance-$d$ surface code requires $\approx d^2$ physical qubits per logical qubit and $d$ rounds of syndrome extraction per logical gate in rough estimates; exact constants depend on implementation and magic-state factories.
\end{example}

\section{Noise, Channels, and Mitigation}
\label{sec:nisq}

\subsection{NISQ definition}

NISQ devices \cite{preskill2018} have $50$--$500$ noisy qubits without full error correction. Useful tasks include variational algorithms, simulation prototypes, and benchmarking---not arbitrary-depth fault-tolerant programs.

\subsection{Mitigation versus correction}

Mitigation (Chapter~\ref{ch:qec}) reduces statistical bias in short circuits; correction suppresses logical error rates exponentially given sufficient redundancy. Long-term utility requires the latter for deep algorithms.

\section{Quantum Advantage and Benchmarks}
\label{sec:advantage}

\subsection{Random circuit sampling}

Random circuit sampling milestones \cite{arute2019} demonstrate quantum behavior difficult to simulate classically for \emph{specific} tasks. This is not universal speedup for everyday workloads.

\subsection{How to read claims}

Always ask: \emph{Which problem? Which metric? Which classical comparison?} Cross-entropy benchmarks, quantum volume, and application-specific demos measure different things.

\begin{example}
Sycamore-style supremacy argues classical simulation of a particular random circuit distribution is infeasible; it does not imply Gmail runs faster on quantum processors.
\end{example}

\section{Quantum Networking and Cryptography}
\label{sec:networking-crypto}

\subsection{Quantum key distribution}

QKD (Chapter~\ref{ch:protocols}) offers information-theoretic security under explicit device and channel trust models. Deployment exists in limited links; scaling requires quantum repeaters still under research.

\subsection{Post-quantum classical cryptography}

Shor's algorithm threatens RSA and elliptic-curve cryptography by efficiently solving hidden subgroup problems underlying those schemes. \emph{Post-quantum cryptography} (PQC) standardizes classical algorithms believed secure against quantum adversaries (lattice-based, hash-based, code-based). PQC deployment proceeds independently of quantum hardware timelines.

\section{Software Ecosystem (2026)}
\label{sec:software-ecosystem}

\subsection{Stack layers}

\begin{itemize}[noitemsep]
\item \textbf{Languages \& IR:} OpenQASM~3, Quil, vendor DSLs.
\item \textbf{SDKs:} Qiskit, Cirq, PennyLane, Braket, Azure Quantum, IBM Quantum Platform \cite{ibm2026}.
\item \textbf{Compilers:} routing, pulse-level optimization, error-aware transpilation.
\item \textbf{Simulators:} state-vector, density matrix, stabilizer, tensor-network.
\end{itemize}

\subsection{Hybrid workflows}

Near-term applications combine classical optimizers with quantum coprocessors (variational quantum eigensolvers, QAOA). HPC integration treats quantum devices as specialized accelerators with queueing and calibration schedules.

\section{Fault Tolerance Roadmaps}
\label{sec:ft-roadmaps}

\subsection{Magic states and distillation}

Non-Clifford gates (especially $T$) dominate resource estimates. \emph{Magic-state distillation} factories supply clean $T$ states to fault-tolerant circuits; their footprint often exceeds the computation proper.

\subsection{Company versus science}

Vendor announcements (e.g.\ processor codenames, qubit counts) are engineering milestones, not proofs of utility. Distinguish peer-reviewed logical-qubit demonstrations \cite{acharya2024} from marketing materials \cite{google2024}.

\section{Resource Estimation and Classical Simulation}
\label{sec:resources}

Exact simulation of $n$ qubits needs $O(2^n)$ complex amplitudes---infeasible beyond $\approx 50$ qubits on classical supercomputers for generic states. Clifford circuits simulate in polynomial time, clarifying why non-Clifford content matters for advantage claims.

\section{Open Questions at the \cutoff{} Cutoff}
\label{sec:open-questions}

\begin{itemize}
\item Fault-tolerant logical qubits at scale with practical space-time overheads.
\item Durable quantum advantage in chemistry, optimization, or machine learning.
\item Integration of quantum accelerators with classical HPC and secure PQC migration.
\item Standardized benchmarks correlating device metrics with application performance.
\end{itemize}

\section*{Further Reading}

See \cite{nielsen2010,preskill2018,ibm2026} and peer-reviewed surveys. Re-verify hardware claims against primary literature before citing in research.

\section*{Exercises}
\begin{enumerate}[label=\arabic*.]
\item Compare two hardware platforms on fidelity, coherence, and connectivity in a table.
\item Explain why surface codes favor 2D-local architectures.
\item Distinguish quantum advantage, quantum supremacy, and fault tolerance in one paragraph each.
\item Why does PQC migration not require a working quantum computer?
\item Estimate qubits needed to factor RSA-2048 using rough literature exponents (order-of-magnitude).
\item Name three software layers in the 2026 quantum stack.
\item What role do magic states play in fault-tolerant $T$ gates?
\item Critique a hypothetical headline: ``1000-qubit computer breaks encryption.''
\end{enumerate}


\appendix
\chapter{Gate Reference}
Standard single-qubit gates:

\[
I=\begin{bmatrix}1&0\\0&1\end{bmatrix},\quad
X=\begin{bmatrix}0&1\\1&0\end{bmatrix},\quad
Y=\begin{bmatrix}0&-i\\i&0\end{bmatrix},\quad
Z=\begin{bmatrix}1&0\\0&-1\end{bmatrix},
\]
\[
H=\frac{1}{\sqrt2}\begin{bmatrix}1&1\\1&-1\end{bmatrix},\quad
S=\begin{bmatrix}1&0\\0&i\end{bmatrix},\quad
T=\begin{bmatrix}1&0\\0&e^{i\pi/4}\end{bmatrix}.
\]

Rotation gates:
\[
R_x(\theta)=e^{-i\theta X/2},\quad
R_y(\theta)=e^{-i\theta Y/2},\quad
R_z(\theta)=e^{-i\theta Z/2}.
\]

Two-qubit:
\[
\mathrm{CNOT}=\begin{bmatrix}1&0&0&0\\0&1&0&0\\0&0&0&1\\0&0&1&0\end{bmatrix}.
\]

Useful identities: $HXH=Z$, $HZH=X$, $H^2=I$, $T^2=S$, $S^2=Z$.


\chapter{Formula Sheet}
Single qubit: $\ket{\psi}=\alpha\ket{0}+\beta\ket{1}$, $|\alpha|^2+|\beta|^2=1$.

Measurement: $P(i)=|\braket{i|\psi}|^2$.

Bloch: $\ket{\psi}=\cos\frac{\theta}{2}\ket{0}+e^{i\phi}\sin\frac{\theta}{2}\ket{1}$.

Unitary: $U^\dagger U=I$.

Bell: $\ket{\Phi^+}=\frac{\ket{00}+\ket{11}}{\sqrt2}$.

Grover: $O_w=I-2\ket{w}\bra{w}$, $D=2\ket{s}\bra{s}-I$.

QFT: $\mathrm{QFT}\ket{x}=\frac{1}{\sqrt N}\sum_k e^{2\pi ixk/N}\ket{k}$.

Shor: $a^r\equiv 1\pmod N$, factors from $\gcd(a^{r/2}\pm 1,N)$.


\chapter{Complexity Reference}
\begin{tabular}{@{}lll@{}}
\toprule
Problem & Classical & Quantum \\
\midrule
Deutsch ($n=1$) & 2 queries (worst) & 1 query \\
Deutsch--Jozsa & $2^{n-1}+1$ & 1 \\
Bernstein--Vazirani & $n$ & 1 \\
Grover search & $O(N)$ & $O(\sqrt N)$ \\
Factoring (Shor) & no known poly-time & poly-time (ideal model) \\
\bottomrule
\end{tabular}

\medskip
\textbf{BQP} is the class of decision problems solvable in polynomial time with bounded error on a quantum computer. BQP is contained in PSPACE; the relationship to NP is unknown.

\medskip
\textbf{Caution:} Different algorithms provide different advantages. Quantum computers are not universally exponentially faster.


\chapter{Glossary}
\begin{description}[style=nextline, leftmargin=0pt, labelwidth=2.5cm]

\item[Amplitude]
Complex coefficient multiplying a basis ket in a superposition; probabilities derive from squared magnitudes.

\item[Ancilla]
Auxiliary qubit used temporarily during computation (e.g.\ oracle workspace, syndrome measurement).

\item[Basis]
Orthonormal set of vectors spanning a Hilbert space; measurement outcomes label basis elements.

\item[Bit]
Classical two-level information unit, $0$ or $1$.

\item[Bloch sphere]
Geometric representation of pure single-qubit states as points on the unit sphere, modulo global phase.

\item[BQP]
Bounded-error quantum polynomial time; decision problems efficiently solvable on a quantum computer.

\item[Circuit depth]
Number of sequential time layers in a quantum circuit, counting parallel gates on disjoint qubits as one layer.

\item[Clifford gate]
Gate generated by $\{H,S,\mathrm{CNOT}\}$; simulable efficiently classically (Gottesman--Knill).

\item[CNOT]
Controlled-NOT gate; flips target when control is $\ket{1}$.

\item[Decoherence]
Loss of quantum coherence through interaction with an uncontrolled environment.

\item[Depolarizing channel]
Noise model applying random $I,X,Y,Z$ Pauli errors with fixed probability.

\item[Ebit]
One unit of shared entanglement, e.g.\ one Bell pair.

\item[Entanglement]
Non-separable quantum state of composite systems; cannot be written as a product of subsystem states.

\item[Error syndrome]
Classical bit pattern from stabilizer measurements indicating likely error location/type.

\item[Fault tolerance]
Computation scheme maintaining arbitrary low logical error rates given physical errors below a threshold.

\item[Gate]
Reversible operation on bits (classical) or unitary operation on qubits (quantum).

\item[Global phase]
Overall complex phase $e^{i\gamma}$ on a state vector; unobservable in standard measurements.

\item[Hadamard gate ($H$)]
Single-qubit gate exchanging $Z$ and $X$ bases: $H\ket{0}=\ket{+}$.

\item[Hilbert space]
Complex vector space with inner product where quantum states live.

\item[Ket ($\ket{\psi}$)]
Dirac notation for a column state vector.

\item[Bra ($\bra{\psi}$)]
Conjugate transpose row vector dual to $\ket{\psi}$.

\item[Logical qubit]
Error-protected encoded qubit implemented across many physical qubits.

\item[Magic state]
Resource state enabling non-Clifford gates fault-tolerantly, especially $T$ gates.

\item[Measurement]
Process yielding classical outcomes with Born-rule probabilities and state update.

\item[Mitigation]
Statistical techniques reducing noise bias without full quantum error correction.

\item[NISQ]
Noisy intermediate-scale quantum; devices too noisy for full fault tolerance yet useful for experiments.

\item[No-cloning theorem]
Unknown arbitrary quantum states cannot be copied by any unitary process.

\item[No-signaling]
Entanglement cannot transmit controllable information faster than light.

\item[Oracle]
Black-box unitary implementing a function in query complexity and algorithms.

\item[Pauli operators]
Single-qubit gates $X,Y,Z$; building blocks for errors and stabilizers.

\item[Phase]
Argument of a complex amplitude; relative phase affects interference.

\item[Physical qubit]
Hardware-level two-level system subject to noise and calibration drift.

\item[Post-quantum cryptography (PQC)]
Classical cryptographic algorithms conjectured secure against quantum adversaries.

\item[Projector]
Operator $P=\ket{m}\bra{m}$ associated with measurement outcome $m$.

\item[Qubit]
Two-level quantum information unit generalizing the classical bit.

\item[Stabilizer]
Pauli operator commuting with a code space; measured to extract syndromes.

\item[Superposition]
Linear combination of basis states with more than one nonzero amplitude.

\item[Surface code]
Leading 2D topological quantum error-correcting code with local checks.

\item[Tensor product]
Composition of subsystem Hilbert spaces: $\mathcal{H}=\mathcal{H}_A\otimes\mathcal{H}_B$.

\item[Threshold theorem]
Result guaranteeing fault tolerance if physical gate error rates fall below a code threshold.

\item[Unitary]
Linear operator $U$ with $U^\dagger U=I$; reversible quantum gate before measurement.

\end{description}


\backmatter
\bibliographystyle{plain}
\bibliography{references}

\end{document}
